\section{Reproducibility data for the molecular suite}
\label{app:repro}

% --- float declaration moved here (placement only): this table is the LAST
%     float of the manuscript and it used to be declared after the one
%     paragraph of this appendix, with no text left behind it.  The
%     \clearpage that main.tex puts before the bibliography therefore flushed
%     it onto a page of its own carrying not one line of text.  Declared
%     ahead of the paragraph it can top the page the appendix starts on, and
%     that page keeps the paragraph.  Not a character of the table or of the
%     paragraph changed, and declaration ORDER is unchanged, so Table XVI is
%     still Table XVI.
\begin{table*}[tp]\centering\small
\caption{\textbf{Reproducibility data for the nineteen-molecule suite} (cc-pVDZ, valence bond-breaking
active spaces). $R_{\rm eq},R_{\rm diss}$ in \AA{}; for the polyatomics ($\mathrm{BeH_2},\mathrm{H_2O},
\mathrm{NH_3},\mathrm{H_4},\mathrm{H_6}$) the representative varied bond length is shown and full
Cartesians are in the ancillary data. $D$: dimension of the $N$-electron (neutral, ground-state) sector,
in which the tabulated $|\mathcal S|$ lives; $|\mathcal S|$: ground-state determinant support ($\ge99\%$
FCI weight); $\chi$: minimal MPS bond dimension across the central spin-orbital cut. The $|\mathcal S|/D$
of Table~\ref{tab:molecules} is a distinct ratio, taken in the $(N{\pm}1)$ spectral sector.}
\label{tab:repro}
\begin{tabular}{llccccccc}\hline\hline
Molecule & CAS$(n_e,n_{\rm o})$ & $R_{\rm eq}$ & $R_{\rm diss}$ & $D$ & $|\mathcal S|^{\rm eq}$ & $|\mathcal S|^{\rm diss}$ & $\chi^{\rm eq}$ & $\chi^{\rm diss}$ \\ \hline
F$_2$ & (1+1,2) & 1.412 & 3.00 & 4 & 1 & 2 & 2 & 2 \\
N$_2$ & (3+3,6) & 1.098 & 2.40 & 400 & 7 & 80 & 7 & 32 \\
H$_2$O & (2+2,4) & 0.958 & 2.30 & 36 & 1 & 13 & 1 & 13 \\
HF & (1+1,2) & 0.917 & 2.30 & 4 & 1 & 4 & 1 & 4 \\
NH$_3$ & (3+3,6) & 1.012 & 2.33 & 400 & 1 & 68 & 4 & 42 \\
CO & (3+3,6) & 1.128 & 2.40 & 400 & 8 & 46 & 14 & 30 \\
H$_6$ & (3+3,6) & 0.900 & 2.00 & 400 & 9 & 91 & 8 & 36 \\
H$_4$ & (2+2,4) & 0.900 & 2.00 & 36 & 2 & 11 & 4 & 8 \\
H$_2$ & (1+1,2) & 0.741 & 2.00 & 4 & 1 & 2 & 1 & 2 \\
BeH$_2$ & (2+2,4) & 1.334 & 2.90 & 36 & 1 & 6 & 1 & 8 \\
C$_2$ & (4+4,6) & 1.243 & 2.60 & 225 & 13 & 6 & 11 & 2 \\
NO & (4+3,6) & 1.151 & 2.50 & 300 & 5 & 25 & 11 & 7 \\
O$_2$ & (5+3,6) & 1.208 & 2.60 & 120 & 11 & 4 & 4 & 2 \\
CN & (4+3,6) & 1.172 & 2.50 & 300 & 27 & 9 & 17 & 7 \\
OH & (3+2,4) & 0.970 & 2.40 & 24 & 1 & 5 & 1 & 4 \\
BH & (2+2,4) & 1.232 & 2.60 & 36 & 5 & 3 & 2 & 2 \\
BeH & (2+1,3) & 1.343 & 2.70 & 9 & 1 & 2 & 1 & 2 \\
LiH & (1+1,2) & 1.596 & 3.20 & 4 & 1 & 2 & 1 & 4 \\
LiF & (1+1,2) & 1.564 & 3.20 & 4 & 1 & 1 & 1 & 1 \\
\hline\hline\end{tabular}
\end{table*}

Table~\ref{tab:repro} lists, for every molecule of Table~\ref{tab:molecules}, the equilibrium and
dissociation bond lengths, the active space, the $N$-electron (neutral, ground-state) sector dimension $D$, the ground-state
determinant support $|\mathcal S|$ (the number of determinants carrying $\ge99\%$ of the FCI weight), and
the minimal matrix-product-state bond dimension $\chi$ at both geometries. These are the quantities that
enter the $\mathcal F_1$--$|\mathcal S|$ and $\chi$--$|\mathcal S|$ correlations of
Sec.~\ref{sec:nogo}, pooled over the $n=38$ equilibrium-and-dissociation points (Spearman
$\chi$--$|\mathcal S|$ $\rho=0.90$, $\mathcal F_1$--$|\mathcal S|$ $\rho=0.86$; Pearson
$\mathcal F_1$--$\log_2|\mathcal S|$ $r=0.80$). The $|\mathcal S|/D$ reported in
Table~\ref{tab:molecules} is a \emph{distinct} quantity---the fraction of the $(N{+}1)$ sector used in
the spectral \emph{reconstruction}, not the ground-state support tabulated here. Every active-space FCI
energy is converged to better than $10^{-12}$~Ha (the largest residual across all $38$ points is
$1.1\times10^{-13}$~Ha), comfortably below the $10^{-5}$~Ha quoted in the main text; the ground-state
energies, full Cartesian geometries, and orbital active-space indices are provided as ancillary files.
